%%%%%%%%%%%%%%%%%%%%%%%%%%%%%%%%%%%%%%%%%
% Plik z wstępem do pracy
% Szablon pracy dyplomowej
% Wydział Informatyki 
% Zachodniopomorski Uniwersytet Technologiczny w Szczecinie
% autor Joanna Kołodziejczyk (jkolodziejczyk@zut.edu.pl)
% Bardzo wczesnym pierwowzorem szablonu był
% The Legrand Orange Book
% Version 5.0 (29/05/2025)
%
% Modifications to LOB assigned by %JK
%%%%%%%%%%%%%%%%%%%%%%%%%%%%%%%%%%%%%%%%%


\chapter*{Wstęp}

Rozwój technologii komputerowych nieustannie przesuwa granice możliwości
interakcji człowieka z maszyną. Tradycyjne urządzenia wejściowe ---
klawiatura i mysz --- wymagają sprawności ruchowej, której nie posiadają
osoby dotknięte chorobami neurologicznymi takimi jak stwardnienie zanikowe
boczne, porażenie czterokończynowe czy zespół zamknięcia. Dla tej grupy
pacjentów interfejsy mózg-komputer (ang.~\textit{Brain-Computer Interface},
BCI) stanowią często jedyną możliwość komunikacji ze światem zewnętrznym
i samodzielnego sterowania urządzeniami. Rosnące zainteresowanie tą
technologią, postęp w miniaturyzacji sprzętu EEG oraz dostępność
bibliotek do przetwarzania sygnałów biomedycznych sprawiają, że budowa
funkcjonalnego systemu BCI stała się możliwa poza środowiskiem czysto
laboratoryjnym.

Spośród paradygmatów stosowanych w systemach BCI, wzrokowe potencjały
wywołane stanu ustalonego (ang.~\textit{Steady-State Visual Evoked
Potentials}, SSVEP) wyróżniają się wysoką skutecznością klasyfikacji,
krótkim czasem uczenia się przez użytkownika oraz brakiem konieczności
indywidualnej kalibracji klasyfikatora. Zjawisko to polega na
synchronizacji aktywności neuronów kory wzrokowej z rytmem migającego
bodźca wizualnego, co generuje charakterystyczne prążki w widmie sygnału
EEG. Umożliwia to jednoznaczne powiązanie intencji użytkownika
z konkretną komendą sterującą na podstawie samego skupienia wzroku.

Niniejsza praca podejmuje problem zaprojektowania i implementacji
interfejsu SSVEP umożliwiającego sterowanie kursorem myszy w czasie
rzeczywistym przy użyciu ogólnodostępnego sprzętu EEG. Celem pracy
jest zbudowanie działającego systemu BCI opartego na płytce OpenBCI
Cyton, weryfikacja jego skuteczności w badaniu z udziałem użytkowników
oraz ocena użyteczności interfejsu. Przyjęto tezę, że system oparty
na analizie korelacji kanonicznej (CCA) jest w stanie poprawnie
klasyfikować komendy sterujące u większości użytkowników bez
przeprowadzania indywidualnej sesji treningowej.

Cel został osiągnięty poprzez kolejne etapy: analizę literatury
dotyczącej metod przetwarzania sygnału SSVEP, projekt architektury
wieloprocesowej systemu, implementację modułów akwizycji, filtracji,
klasyfikacji i stymulacji wizualnej, a następnie przeprowadzenie
badania z trzema uczestnikami obejmującego zadanie sterowania
kursorem oraz ocenę skuteczności klasyfikacji komend.

Praca składa się z czterech rozdziałów. Rozdział~pierwszy wprowadza
pojęcie interfejsu mózg-komputer, omawia metody rejestracji aktywności
mózgu ze szczególnym uwzględnieniem elektroencefalografii oraz opisuje
rytmy EEG stanowiące podstawę działania systemów BCI. Rozdział~drugi
przedstawia zjawisko SSVEP, strukturę interfejsu opartego na tym
paradygmacie, metody prezentacji bodźców wizualnych oraz algorytmy
przetwarzania sygnału --- analizę widmową i analizę korelacji
kanonicznej. Rozdział~trzeci zawiera opis projektu i implementacji
systemu: założenia projektowe, architekturę modułową, dobór częstotliwości
stymulacji, implementację klasyfikatora CCA oraz mechanizm bufora
decyzyjnego. Rozdział~czwarty przedstawia wyniki badania przeprowadzonego
z udziałem trzech uczestników wraz z omówieniem uzyskanych wyników
i wnioskami dotyczącymi dalszego rozwoju systemu.